\chapter{RITI SACRIFICALI}\label{riti-sacrificali}

\hfill\break

Jason, ossessionato dall'imminente arrivo di E.D. Lawton, aveva
trascurato di menzionare il fatto che alla Perielio era atteso anche un
altro ospite: Preston Lomax, l'allora Vicepresidente degli Stati Uniti e
favorito nelle successive elezioni.

Ai cancelli la sicurezza era serrata e c'era un elicottero sulla
piattaforma centrale, in cima all'edificio della Perielio. Riconobbi
tutti quei protocolli da codice rosso perché nell'ultimo mese c'era
stata una serie di visite del Presidente Garland. La guardia
all'ingresso principale, quella che mi chiamava ``Doc'' e a cui una
volta al mese monitoravo il livello del colesterolo, mi informò che si
trattava di Lomax.

Avevo appena superato la porta della clinica (in assenza di Molly si
occupava della reception una sostituta di nome Lucinda) quando ricevetti
un messaggio al cercapersone che mi fece deviare verso l'ufficio di
Jason nell'ala dirigenziale. Dopo aver attraversato quattro perimetri di
sicurezza mi ritrovai da solo con lui. Temevo che volesse chiedermi
altri farmaci, ma la terapia della sera prima aveva migliorato le sue
condizioni in maniera convincente seppure solo temporanea. Si alzò e mi
venne incontro, mostrandomi la mano tesa che non tremava: «Volevo
ringraziarti per questa, Ty.»

«Figurati, ma devo ripetermi\ldots{} non posso garantirti nulla.»

«Recepito. Basta che stia bene per oggi. E.D. arriva a mezzogiorno.»

«Tralasciando il Vicepresidente.»

«Lomax è qui da stamattina alle sette. È un tipo mattiniero. È rimasto
un paio di ore a colloquio con il nostro ospite marziano e tra poco mi
toccherà fargli da guida per un giro di cortesia. E a proposito di Wun,
gli piacerebbe vederti, se hai qualche minuto libero.»

«Sempre che gli affari nazionali non lo stiano tenendo occupato.» Era
molto probabile che la settimana successiva Lomax avrebbe ottenuto il
voto della nazione - una passeggiata, stando ai sondaggi. Era da molto
prima dell'arrivo di Wun che Jase si teneva buono Lomax e questi era
affascinato da Wun. «Anche tuo padre si unirà al giro?»

«Solo perché non c'è un modo educato per tenerlo fuori.»

«Prevedi qualche problema?»

«Prevedo molti problemi.»

«Fisicamente, comunque, stai bene?»

«Mi sento bene, ma il dottore sei tu. Mi servono solo un paio d'ore in
più, Tyler. Immagino di farcela, no?»

La frequenza dei battiti era un po' alta - com'era prevedibile - ma i
sintomi della SMA erano ben contenuti. A guardarlo, non si poteva dire
se le medicine gli avessero lasciato agitazione o confusione. In effetti
era di una calma quasi radiosa, chiuso in qualche stanza lucida e
imperturbabile in fondo alla sua mente.

E così andai a trovare Wun Ngo Wen. Wun non era nel suo alloggio; si era
spostato nella piccola mensa dei dirigenti, transennata e circondata da
colossi con gli auricolari dal cavetto attorcigliato nascosto dietro le
orecchie. Quando passai di fianco al carrello, sollevò lo sguardo e
mandò via con un cenno i cloni della sicurezza che si erano avvicinati
per fermarmi.

Mi sedetti di fronte a lui a un tavolo con la superficie in vetro.
Piluccò un insipido trancio di salmone con una forchetta della mensa e
sorrise sereno. Mi incurvai sulla sedia per essere al suo livello. Gli
avrebbe fatto comodo un rialzo.

Ma il cibo gli faceva bene. Aveva messo su qualche chilo nel periodo
passato alla Perielio, pensai. Il suo abito, fatto su misura un paio di
mesi prima, gli andava stretto sul ventre. Aveva evitato di abbottonarsi
il panciotto abbinato. Anche le guance erano più piene malgrado fossero
rugose come sempre, la pelle scura lievemente afflosciata.

«So che hai avuto una visita» dissi.

Wun annuì. «Ma non è stata la prima volta. Ho incontrato il Presidente
Garland a Washington in diverse occasioni e ho incontrato due volte il
Vicepresidente Lomax. Si prevede che le elezioni lo porteranno al
potere, così dicono.»

«Non perché sia particolarmente amato.»

«Non sono in grado di giudicarlo come candidato, ma fa domande davvero
interessanti.»

Quell'approvazione mi fece sentire un po' protettivo. «Sono certo che sa
essere affabile quando vuole. E durante il suo mandato ha fatto un
lavoro decente, ma per buona parte della sua carriera è stato l'uomo più
odiato di Capitol Hill. Capogruppo in tre diverse amministrazioni. Non
gli sfugge quasi niente.»

Wun fece un gran sorriso. «Pensi che io sia ingenuo, Tyler? Temi che il
Vicepresidente Lomax si approfitti di me?»

«Ingenuo no, non proprio\ldots»

«Sono un principiante, in effetti. Mi sfuggono le sottili sfumature
politiche, ma ho diversi anni più di Preston Lomax e io stesso ho
ricoperto una carica pubblica.»

«Davvero?»

«Per tre anni,» disse con percettibile orgoglio, «sono stato Ispettore
dell'agricoltura nel Distretto dei Venti Ghiacciati.»

«Ah.»

«L'organo governativo di gran parte del delta del Kirioloj. Non era la
Presidenza degli Stati Uniti d'America. Non ci sono armi atomiche a
disposizione dell'Ispettorato dell'agricoltura. Ma denunciai un
funzionario locale corrotto che falsificava i rapporti dei raccolti
alterandone il peso e vendeva il suo margine al mercato delle
eccedenze.»

«Un sistema di tangenti?»

«Se si dice così\ldots»

«Quindi le Cinque Repubbliche non sono esenti da corruzione?»

Wun batté le palpebre, un gesto che si propagò sulla geografia convoluta
del suo volto. «No, e come potrebbero esserlo? Perché a tanti terrestri
sembra strana questa cosa? Se fossi arrivato da qualche altro paese
della Terra - Francia, Cina, Texas - nessuno si sarebbe stupito di
sentir parlare di corruzione, disonestà o furto.»

«Credo di no. Ma non è la stessa cosa.»

«No? Ma tu lavori qui alla Perielio. Devi aver conosciuto qualcuno della
generazione fondatrice - per quanto mi sembri ancora strano questo
concetto - le donne e gli uomini di cui noi marziani siamo i lontani
discendenti. Erano persone così perfette da aspettarti che la loro
progenie fosse senza peccato?»

«No, ma\ldots»

«Eppure l'equivoco è quasi universale. Anche quei libri che mi hai dato,
scritti prima dello Spin\ldots»

«Li hai letti?»

«Sì, con avidità. Mi sono piaciuti. Grazie. Ma anche in quei romanzi, i
marziani\ldots» si sforzò di cercare le parole giuste.

«Sì, alcuni sembrano quasi dei santi\ldots» gli completai la frase.

«Distaccati, saggi, in apparenza fragili, ma in realtà molto potenti.
Gli Antichi. Ma per noi, Tyler, siete \emph{voi} gli Antichi. La specie
più vecchia, il pianeta antico. Trovo tutto questo molto ironico.»

Ci riflettei su. «Anche il romanzo di HG. Wells\ldots»

«I suoi marziani si vedono a malapena. Sono cattivi in maniera astratta,
indifferente. Non saggi, bensì intelligenti. Ma diavoli e angeli sono
fratelli e sorelle, se ho interpretato bene le vostre tradizioni
popolari.»

«Ma le storie più contemporanee\ldots»

«Quelle erano davvero interessanti e i protagonisti almeno erano umani.
Ma il piacere più puro di quelle storie sta nei paesaggi, non sei
d'accordo? E ciò nonostante sono paesaggi \emph{trasformativi}. Un
destino dietro ogni duna.»

«E naturalmente Bradbury\ldots»

«Il suo Marte non è Marte, ma il suo Ohio ci si avvicina molto.»

«Ho capito quello che vuoi dire. Siete semplicemente persone e Marte non
è il paradiso. D'accordo, ma questo non significa che Lomax non proverà
a usarti per i suoi scopi politici.»

«E io ti assicuro che sono perfettamente consapevole di questa
possibilità. O, meglio, di questa \emph{certezza}. È ovvio che verrò
usato per un vantaggio politico, ma è questo il potere che ho: concedere
o negare il mio appoggio. Collaborare o ribellarmi. Il potere di dire la
parola giusta.» Sorrise di nuovo. Aveva denti perfetti e uniformi e di
un bianco splendente. «O di non farlo.»

«Quindi \emph{cosa} vuoi guadagnarci?»

Mi mostrò i palmi, un gesto terrestre ma anche marziano. «Niente. Sono
un santo marziano, no? Ma sarebbe gratificante vedere realizzato il
lancio dei replicatori.»

«Solo per il bene della conoscenza?»

«Sì, lo \emph{confesso}, anche se è una motivazione \emph{da santo}.
Conoscere almeno qualcosa sullo Spin\ldots»

«E sfidare gli Ipotetici?»

Batté di nuovo le palpebre. «Spero con tutto il cuore che gli Ipotetici,
chiunque o qualunque cosa siano, non percepiscano quello che faremo come
una \emph{sfida}.»

«Ma se dovesse accadere\ldots»

«Perché dovrebbe?»

«Ma se dovesse accadere, crederanno che la sfida provenga dalla Terra,
non da Marte.»

Wun Ngo Wen batté le palpebre diverse altre volte. Poi il sorriso
riapparve lentamente: benevolo, di approvazione. «Sei davvero cinico,
tu, dottor Dupree.»

«Proprio da non-marziano.»

«Davvero.»

«Invece Preston Lomax crede che tu sia un angelo?»

«Solo lui può rispondere a questa domanda. L'ultima cosa che mi ha
detto\ldots» E qui Wun abbandonò la sua dizione di Oxford per imitare in
modo perfetto Preston Lomax, brusco e gelido come una spiaggia
d'inverno: «È un privilegio parlare con lei, Ambasciatore Wen. Dice ciò
che pensa in modo chiaro. Molto piacevole per un veterano di D.C. come
me.»

L'imitazione era sorprendente, per uno che parlava inglese da poco più
di un anno. Glielo dissi.

«Sono uno studioso» rispose. «Leggo l'inglese da quando ero piccolo.
Parlarlo è un'altra cosa. Ma sono portato per le lingue. È uno dei
motivi per cui sono qui. Tyler, posso chiederti un altro favore? Saresti
disposto a portarmi altri romanzi?»

«Temo di aver finito le storie marziane.»

«Non su Marte. Qualsiasi genere di romanzo. Tutto, tutto quello che
consideri importante, tutto quello che per te conta o che ti ha dato un
po' di piacere.»

«Ci saranno un sacco di professori che sarebbero felici di prepararti
una lista di letture.»

«Certo che ci sono, ma io lo sto chiedendo a te.»

«Non sono uno studioso. Mi piace leggere, ma quello che mi capita e
soprattutto contemporanei.»

«Tanto meglio. Sono solo più spesso di quanto tu creda. Il mio alloggio
è comodo, ma non posso allontanarmi senza una minuziosa pianificazione.
Non posso andare a mangiare fuori, non posso vedere un film o
frequentare un circolo. Potrei chiedere dei libri alle mie guardie del
corpo, ma l'ultima cosa che voglio è un'opera di narrativa approvata da
una commissione. Un libro sincero, invece, è quasi come un buon amico.»

Era la cosa più simile a una lamentela a cui Wun fosse giunto riguardo
alla sua posizione alla Perielio, alla sua posizione sulla Terra. Nelle
ore in cui era sveglio si sentiva abbastanza felice, disse, troppo
occupato per provare nostalgia e ancora euforico per la stranezza di ciò
che per lui sarebbe stato sempre un mondo alieno. Ma di notte, sul punto
di dormire, a volte immaginava di camminare sulla riva di un lago
marziano, di osservare uccelli acquatici riunirsi in stormi e
volteggiare sulle onde, e nella sua mente era sempre un pomeriggio
nebbioso, colorato da strisce dell'antica sabbia che continuava a
sollevarsi dai deserti di Noachis per dipingere il cielo. In questo
sogno o visione era da solo, disse, ma sapeva che c'erano altri ad
aspettarlo dietro la vicina curva della costa rocciosa. Potevano essere
amici o sconosciuti, poteva essere anche la sua famiglia perduta; sapeva
soltanto che lo avrebbero accolto, toccato, stretto a loro, abbracciato.
Ma era solo un sogno.

«Quando leggo,» mi disse, «sento l'eco di quelle voci.»

Gli promisi che gli avrei portato dei libri, però in quel momento ci fu
altro da fare. Nacque un improvviso fermento nel cordone di sicurezza
vicino alla porta della mensa. Uno dei funzionari si avvicinò e gli
disse: «Chiedono di voi al piano di sopra.»

Wun smise di mangiare e cominciò a scendere con difficoltà dalla sedia.
Gli dissi che ci saremmo visti più tardi.

Il funzionario si rivolse a me. «Anche lei» disse. «Chiedono di
entrambi.»

ooo

La sicurezza ci sollecitò a entrare in una sala riunioni adiacente
all'ufficio di Jason, dove Jase e un manipolo di capidivisione della
Perielio si stava confrontando con una delegazione che includeva E.D.
Lawtone il probabile futuro Presidente, Preston Lomax. Nessuno sembrava
felice.

Guardai verso E.D. Lawton, che non vedevo dal funerale di mia madre. La
sua estrema magrezza cominciava a sembrare quasi patologica, come se si
fosse svuotato di qualcosa di vitale. Polsini bianchi inamidati, polsi
ossuti abbronzati. I capelli erano radi, flosci e pettinati a casaccio,
ma gli occhi ancora vigili. Gli occhi di E.D erano sempre vividi quando
era arrabbiato.

Preston Lomax, invece, sembrava solo impaziente. Era venuto alla
Perielio per farsi fotografare con Wun (foto per la stampa dopo
l'annuncio ufficiale della Casa Bianca) e per conferire sulla strategia
dei replicatori che programmava di appoggiare. E.D. era lì grazie alla
sua reputazione. Era stato lui, in persona, a parlare durante il tour
preelettorale del Vicepresidente, e da allora evidentemente non aveva
più smesso.

Durante la visita, durata un'ora, E.D. aveva contestato, messo in
dubbio, deriso o guardato con sospetto allarmistico praticamente ogni
dichiarazione rilasciata dai capidivisione di Jason, soprattutto quando
la delegazione passò serpeggiando davanti ai nuovi laboratori con le
incubatrici. Ma (secondo Jenna Wylie, la caposquadra della crionica, che
me lo spiegò in seguito) Jason aveva risposto a ogni attacco di suo
padre confutandolo a modo suo, in maniera paziente e forse già ben
preparata. Questo aveva portato E.D. a un alto livello di indignazione,
che lo fece apparire, secondo Jenna, ``come un re Lear impazzito che
farnetica di perfidi marziani.''

La battaglia era ancora in corso quando io e Wun entrammo. E.D., chino
sul tavolo della conferenza, diceva: «Morale della favola, non ha
precedenti, non è sperimentato e include una tecnologia che non
comprendiamo e non possiamo controllare.»

E Jason sorrideva come chi è troppo educato per mettere in imbarazzo un
anziano rispettato ma irascibile. «Ovvio, niente di ciò che facciamo è
privo di rischi. Ma\ldots»

Ma eccoci là. Alcuni dei presenti non avevano mai visto di persona Wun
e, nel riconoscerlo, lo fissarono come pecore spaventate. Lomax si
schiarì la voce. «Scusatemi, ora devo scambiare qualche parola con Jason
e i nostri nuovi arrivati\ldots{} in privato, vi dispiace? Giusto un
attimo.»

Allora la folla obbediente uscì in fila, compreso E.D., che comunque non
aveva l'aria di chi è stato congedato e al contrario sembrava
trionfante.

Porte chiuse. Il silenzio imbottito della sala riunioni si posò attorno
a noi come neve fresca. Lomax, che non ci aveva ancora salutati, si
rivolse a Jason. «So che mi avevi detto che avremmo ricevuto delle
critiche. Così, però\ldots»

«Non è una cosa facile da gestire. Lo capisco.»

«Non mi piace che E.D. pisci dentro alla tenda da fuori. È sconveniente.
Ma in realtà non può nuocerci, supponendo\ldots»

«Supponendo che non ci siano basi per le sue accuse. Te l'assicuro, non
ci sono.»

«Pensi che si sia rimbambito.»

«Non mi spingerei così oltre. Se penso che il suo parere sia diventato
discutibile? Sì, lo penso.»

«Sai che le stesse accuse arrivano da una parte e dall'altra.»

Non ero mai stato tanto vicino, e non lo sarei più stato, a un
Presidente in carica. Lomax non era ancora stato eletto, ma tra lui e
l'incarico restavano solo formalità. Come Vicepresidente, Lomax era
sempre parso un po' severo, un po' meditabondo, una roccia del Maine
rispetto al Texas esuberante di Garland, la presenza ideale per un
funerale di Stato. Durante la campagna aveva imparato a sorridere più
spesso ma lo sforzo non era mai del tutto convincente; i vignettisti
politici ne accentuavano inevitabilmente il cipiglio, mentre si mordeva
il labbro inferiore come se stesse trattenendo una maledizione, gli
occhi freddi come un inverno a Cape Cod.

«Da una parte e dall'altra. Parli delle insinuazioni di E.D. sulla mia
salute.»

Lomax sospirò. «Francamente, il parere di tuo padre sull'aspetto pratico
del progetto dei replicatori non ha molto peso. È un punto di vista
minoritario e probabilmente resterà tale. Però sì, devo ammettere che le
accuse che ha fatto oggi sono un po' fastidiose.» Si voltò verso di me.
«Ecco perché lei è qui, dottor Dupree.»

In quel momento Jason rivolse l'attenzione su di me, la sua voce era
cauta, prudentemente neutrale. «Pare che E.D. faccia delle affermazioni
piuttosto avventate. Dice che soffro di, com'era, una malattia cerebrale
aggressiva\ldots?»

«Un deterioramento neurologico incurabile,» precisò Lomax, «che
interferisce con la capacità di Jason di sovrintendere alle operazioni
qui alla Perielio. Cosa ne dice lei, dottor Dupree?»

«Dico che Jason può parlare per se stesso.»

«L'ho già fatto» ribatté Jase. «Al Vicepresidente Lomax ho spiegato
tutto della mia sclerosi multipla.»

Di cui in realtà non soffriva. Era un segnale. Mi schiarii la gola. «La
sclerosi multipla non è del tutto curabile, ma si può tenere facilmente
sotto controllo. Oggigiorno un paziente di SM può avere un'aspettativa
di vita lunga e produttiva quanto chiunque altro. Forse Jase è stato
riluttante a parlarne, e questo è un suo diritto, ma non c'è niente di
cui vergognarsi, nell'avere la SM.»

Jase mi lanciò un'occhiata insistente che non riuscii a interpretare.
Lomax disse secco: «Grazie, la ringrazio per l'informazione. A
proposito, conosce per caso un certo dottor Malmstein? David Malmstein?»
domandò, seguito da un silenzio che si spalancò come le ganasce di una
trappola d'acciaio.

«Sì» dissi, forse un attimo troppo tardi.

«Questo dottor Malmstein è un neurologo, vero?»

«Sì, lo è.»

«Lo ha consultato in passato?»

«Mi consulto con molti specialisti. Fa parte del mio lavoro di medico.»

«Perché, secondo E.D., ha chiamato questo Malmstein per, ehm, il grave
disturbo neurologico di Jason.»

Questo spiegava l'occhiata gelida che Jase mi stava lanciando. Qualcuno
ne aveva parlato con E.D. Qualcuno con cui aveva un rapporto stretto. Ma
non ero stato io.

Cercai di non pensare a chi potesse essere stato. «Farei lo stesso per
qualsiasi paziente con un'eventuale diagnosi di SM. La clinica che
dirigo qui alla Perielio è buona, ma non abbiamo quelle apparecchiature
diagnostiche a cui può accedere Malmstein in ospedale.»

Credo che Lomax si accorse della mia non risposta, ma rilanciò la palla
a Jase: «Il dottor Dupree dice la verità?»

«Certo che sì.»

«Ti fidi di lui?»

«È il mio medico personale. Certo che mi fido di lui.»

«Perché, senza offesa, ti auguro ogni bene, ma non me ne frega veramente
un cazzo dei tuoi problemi medici. Mi preoccupa solo che tu possa darci
il sostegno necessario e seguire questo progetto fino alla fine. Puoi
farlo?»

«Finché siamo finanziati, sì, signore, sarò qui.»

«E lei che ne pensa, Ambasciatore Wen? Questo le solleva qualche timore?
Qualche preoccupazione o domanda sul futuro della Perielio?»

Wun contrasse le labbra in tre quarti di sorriso marziano. «Davvero
nessuna preoccupazione. Mi fido completamente di Jason Lawton. E mi fido
anche del dottor Dupree. È anche il mio medico personale.»

La cosa costrinse sia me che Jason a trattenere il nostro stupore, ma
siglò l'accordo con Lomax. Scrollò le spalle. «Va bene. Mi scuso per
averne parlato. Jason, spero che tu rimanga in buona salute e che non ti
sia sentito offeso dal tono delle domande, ma data la situazione con
E.D. mi sentivo in dovere di chiedertelo.»

«Capisco» disse Jase. «Riguardo a E.D\ldots»

«Non preoccuparti di tuo padre.»

«Mi dispiacerebbe vederlo umiliato.»

«Verrà messo da parte senza scalpore. Credo sia scontato. Se dovesse
insistere nel fare rivelazioni in pubblico\ldots» Scrollò le spalle. «In
quel caso temo che sarà la \emph{sua} capacità mentale a essere messa in
dubbio.»

«Certo,» disse Jason, «speriamo tutti che non si renda necessario.»

ooo

Passai l'ora successiva in clinica. Quella mattina Molly non si era
presentata e Lucinda aveva fatto tutte le prenotazioni. La ringraziai e
le dissi di prendersi il resto della giornata libera. Pensai di fare un
paio di telefonate, ma non volevo che passassero dal sistema della
Perielio.

Aspettai finché non vidi decollare l'elicottero di Lomax e andar via il
suo corteo maestoso dai cancelli principali; poi riordinai la scrivania
e cercai di pensare a cosa volevo fare. Notai che le mani mi tremavano
un po'. Non era SM. Rabbia, forse. Sdegno. Dolore. Volevo
diagnosticarlo, non provarlo. Volevo relegarlo nell'indice del
\emph{Manuale diagnostico e statistico}.

Avevo superato la reception quando Jason entrò dalla porta.

Disse: «Voglio ringraziarti per avermi dato man forte. Immagino quindi
che non sia stato tu a dire a E.D. di Malmstein.»

«Non lo farei, Jase.»

«Ti credo, ma qualcuno l'ha fatto. E questo rappresenta un problema.
Perché\ldots{} quante persone sanno che vedo un neurologo?»

«Tu, io, Malmstein, chiunque lavori nel suo ufficio\ldots»

«Malmstein non sapeva che E.D. fosse a caccia di pettegolezzi e nemmeno
il suo personale. E.D. deve averlo scoperto da una fonte più vicina. Se
non da me o da te\ldots»

Molly. Non c'era bisogno di dirlo.

«Non possiamo incolparla senza prove.»

«Parla per te. Sei tu quello che ci va a letto. Hai conservato i
risultati dei miei incontri con Malmstein?»

«Non qui in ufficio.»

«A casa?»

«Sì.»

«Gliel'hai mostrati?»

«Certo che no.»

«Ma avrebbe potuto accedervi a tua insaputa.»

«Suppongo di sì.» \emph{Sì}.

«E oggi non è qui a rispondere a eventuali domande. Ha telefonato per
avvisare che stava male?»

Scrollai le spalle. «Non ha proprio chiamato. Lucinda ha cercato di
contattarla, ma non risponde al telefono.»

Jason sospirò. «Non te ne faccio una colpa, ma devi ammettere che hai
fatto molte scelte discutibili in questa storia.»

«Me ne occuperò io» dissi.

«So che sei arrabbiato. Ferito e arrabbiato. Non voglio che tu esca di
qui e faccia qualcosa che peggiorerà le cose. Ma voglio comunque che tu
rifletta sulla tua posizione in questo progetto. Da che parte stai.»

«Lo so, da che parte sto» risposi.

ooo

Cercai di contattare Molly dall'auto, ma continuava a non rispondere. Mi
diressi verso casa sua. Era una giornata calda. Il complesso di case
basse in stucco dove viveva era avvolto dalla foschia degli irrigatori
da giardino. L'odore di funghi della terra bagnata permeava la macchina.

Stavo svoltando verso il parcheggio dei visitatori quando intravidi Moll
che sistemava delle scatole nella parte posteriore di un rimorchio a
noleggio, bianco e malconcio, della U-Haul, agganciato al paraurti
posteriore della sua Ford di tre anni. Mi fermai davanti a lei. Mi notò
ed esclamò qualcosa che non riuscii a sentire, ma che assomigliava molto
a ``Oh, merda!'', però quando scesi dalla macchina rimase dov'era.

«Non puoi parcheggiare là» disse. «Stai bloccando l'uscita.»

«Vai da qualche parte?»

Molly caricò una scatola di cartone con l'etichetta {PIATTI} sul piano
mandorlato del rimorchio. «Secondo te?»

Indossava dei pantaloni sportivi marrone chiaro, una camicia di jeans e
un fazzoletto legato sui capelli. Mi avvicinai e lei indietreggiò di
circa tre passi, chiaramente spaventata.

«Non voglio farti del male» la rassicurai.

«Allora che vuoi?»

«Voglio sapere chi ti ha assunto.»

«Non so di cosa parli.»

«Hai trattato con E.D. in persona o ha usato un intermediario?»

«Merda» disse, valutando la distanza tra lei e la portiera dell'auto.
«Lasciami andare, Tyler. Cosa vuoi da me? Che senso ha questa cosa?»

«Sei andata da lui e gli hai fatto un'offerta o ti ha chiamato lui per
primo? E quando è iniziato tutto, Moll? Mi scopavi per rubare
informazioni fin dall'inizio oppure mi hai venduto a un certo punto dopo
il primo appuntamento?»

«Va' al diavolo.»

«Quanto ti ha pagato? Vorrei sapere quanto valgo.»

«Vaffanculo. Che importa, poi? Non è\ldots»

«Non dirmi che non si tratta di soldi. C'è di mezzo qualche
\emph{principio}?»

«È il \emph{denaro} il principio.» Si pulì le mani sporche di polvere
sui pantaloni, era un po' meno spaventata adesso, un po' più sprezzante.

«Cos'è che vuoi comprare, Moll?»

«Cosa voglio \emph{comprare}? L'unica cosa importante che si
\emph{possa} comprare. Una morte migliore. Una morte migliore, più
pulita. Una di queste mattine sorgerà il Sole e non \emph{smetterà} di
salire finché tutto il cielo del cazzo non sarà bruciato. Allora
scusami, ma io nell'attesa voglio vivere in qualche bel posto. Da sola.
Qualche posto che posso rendere il più confortevole possibile. E quando
arriverà quell'ultima mattina voglio dei farmaci costosi che mi facciano
volare via. Voglio andare a dormire prima che comincino le urla.
Davvero, Tyler. È tutto ciò che voglio, è l'unica cosa in questo mondo
che voglio veramente, e ti ringrazio, grazie per averla resa possibile.»
Era accigliata e arrabbiata, ma le scappò una lacrima che le scivolò
sulla guancia. «Per favore, sposta la macchina.»

«Una bella casa e una bottiglia di pillole? È questo il tuo prezzo?»

«Non c'è nessuno che si prenda cura di me, a parte me.»

«Può sembrare patetico, ma pensavo che potessimo prenderci cura l'uno
dell'altra.»

«Questo significherebbe fidarmi di te. Senza offesa, ma\ldots{}
guardati. Scivoli sulla vita come se aspettassi una risposta o un
salvatore o come se fossi in perenne attesa.»

«Cerco di essere ragionevole, Moll.»

«Ah, non ne dubito. Se la ragionevolezza fosse un coltello starei
sanguinando. Povero ragionevole Tyler. Ma anch'io l'ho capito. È
vendetta, vero? Tutta quella candida santità che indossi come un corredo
di vestiti. È la tua vendetta sul mondo per averti deluso. Il mondo non
ti ha dato ciò che vuoi e non dai nulla indietro se non compassione e
aspirine.»

«Molly\ldots»

«E non osare dire che mi ami, perché so che non è vero. Non conosci la
differenza tra \emph{l'essere} innamorato e \emph{comportarti} come se
lo fossi. È bello che tu mi abbia scelta, ma avresti potuto scegliere
chiunque, e credimi, Tyler, sarebbe stato altrettanto deludente, in un
modo o nell'altro.»

Mi voltai e tornai verso la macchina un po' malfermo, meno scioccato dal
tradimento che dallo scopo, l'intimità cancellata come micro azioni in
un crollo del mercato. Poi mi voltai di nuovo. «E tu, Moll? So che sei
stata pagata per le informazioni, ma è soprattutto per questo che mi hai
scopato?»

«Ti ho \emph{scopato} perché ero \emph{sola}.»

«Sei sola adesso?»

«Non ho mai smesso di esserlo.»

Me ne andai.